\UnnumberedChapter{LỜI MỞ ĐẦU}
Trong các hoạt động chữa cháy và cứu nạn cứu hộ, môi trường khói độc dày đặc và cấu trúc công trình phức tạp luôn đe dọa trực tiếp đến an toàn của lực lượng tác chiến. Khi đó, người chỉ huy ở bên ngoài rất cần theo dõi liên tục trạng thái của từng chiến sĩ trong điều kiện tầm nhìn hạn chế và liên lạc bộ đàm dễ bị gián đoạn. Sự kết hợp giữa công nghệ cảm biến đeo thông minh, thuật toán học máy tại biên và nền tảng IoT mở ra một hướng tiếp cận khả thi để giải quyết bài toán này. Đặc biệt, kỹ thuật định vị tương đối trong nhà đóng vai trò thiết yếu tại những không gian khép kín nơi tín hiệu định vị vệ tinh hoàn toàn bị suy hao.

Xuất phát từ nhu cầu thực tiễn đó, đề tài tập trung thiết kế và xây dựng hệ thống giám sát đa cảm biến theo thời gian thực. Cấu trúc phần cứng tại hiện trường bao gồm cụm thiết bị đeo ở thắt lưng tích hợp vi điều khiển ESP32, cảm biến quán tính GY-85 và khí áp kế BMP580, kết hợp cùng nút quan trắc môi trường gắn trên mũ bảo hộ. Các nút trên cơ thể giao tiếp cục bộ qua BLE và truyền dữ liệu ra ngoài bằng sóng vô tuyến RF 433 MHz. Tại đây, thiết bị đeo trực tiếp thu thập tín hiệu chuyển động để nhận dạng tư thế hoạt động và ước lượng quỹ đạo di chuyển bằng phương pháp định vị quán tính PDR. Tiếp theo, trạm chuyển tiếp Gateway sử dụng máy tính nhúng Raspberry Pi 3 Model B kết hợp mô-đun SIM 4G sẽ tập trung dữ liệu và đồng bộ lên máy chủ MQTT qua mạng di động.

Nội dung báo cáo trình bày hệ thống từ cơ sở lý thuyết, phân tích yêu cầu kỹ thuật đến quy trình thiết kế phần cứng, xây dựng thuật toán tại biên và tổ chức giao thức truyền thông. Bên cạnh đó, hệ thống phần mềm trên nền tảng website và ứng dụng di động Android cũng được phát triển để hiển thị trực quan trạng thái vận động, vị trí tương đối và thông số môi trường. Phần đánh giá kết quả tập trung phân tích tính khả thi của khối thiết bị đeo, mô hình nhận dạng hành động, thuật toán PDR và trạm Gateway.

Nhóm thực hiện xin bày tỏ lòng biết ơn sâu sắc tới TS. Nguyễn Đắc Cử đã tận tình hướng dẫn, định hướng phương pháp luận và đóng góp nhiều ý kiến chuyên môn quý báu trong suốt quá trình nghiên cứu và hoàn thiện đề tài.